% ============================================================================
% French UI and structural strings for the One Math Book series.
% Loaded by styles/onemath.sty when \booklang is fr.
% ============================================================================

% Theorem / statement environment titles
\newcommand{\omnameDefinition}{Définition}
\newcommand{\omnameTheorem}{Théorème}
\newcommand{\omnameProposition}{Proposition}
\newcommand{\omnameLemma}{Lemme}
\newcommand{\omnameCorollary}{Corollaire}
\newcommand{\omnameMethod}{Méthode}
\newcommand{\omnameExample}{Exemple}
\newcommand{\omnameNotation}{Notation}
\newcommand{\omnameRemark}{Remarque}
\newcommand{\omnameExercise}{Exercice}
\newcommand{\omnameExercises}{Exercices}
\newcommand{\omnameProblem}{Problème}
\newcommand{\omnameProblems}{Problèmes}
\newcommand{\omnameProof}{Démonstration}

% Admitted results and solutions appendix headers
\newcommand{\omadmittedtext}{Admis à ce niveau.}
\newcommand{\omsolutionof}{Solution de}
\newcommand{\ompage}{page}
\newcommand{\omnameGotoSolution}{Solution p.}

% Misc text macros
\newcommand{\st}{\text{ tel que }}

% Cover / title page
\newcommand{\bookmaintitle}{One Math Book}
\newcommand{\booktagline}{Le Livre de Maths pour les gouverner tous}
\newcommand{\bookdescription}{Les mathématiques de la maternelle à la fin de la licence}
\newcommand{\bookcontributors}{Les contributeurs du One Math Book}
\newcommand{\bookversionlabel}{Version}

% Colophon
\newcommand{\omcolophonsource}{Code source, signalements et contributions~:}

% Solutions appendix part title
\newcommand{\omsolutionspart}{Solutions des exercices}

% Part titles (Années 1–12; parallel to English Grade N)
\@namedef{omstr@part.grade1}{Année 1 (6--7 ans)}
\@namedef{omstr@part.grade2}{Année 2 (7--8 ans)}
\@namedef{omstr@part.grade3}{Année 3 (8--9 ans)}
\@namedef{omstr@part.grade4}{Année 4 (9--10 ans)}
\@namedef{omstr@part.grade5}{Année 5 (10--11 ans)}
\@namedef{omstr@part.grade6}{Année 6 (11--12 ans)}
\@namedef{omstr@part.grade7}{Année 7 (12--13 ans)}
\@namedef{omstr@part.grade8}{Année 8 (13--14 ans)}
\@namedef{omstr@part.grade9}{Année 9 (14--15 ans)}
\@namedef{omstr@part.grade10}{Année 10 (15--16 ans)}
\@namedef{omstr@part.grade11}{Année 11 (16--17 ans)}
\@namedef{omstr@part.grade12}{Année 12 (17--18 ans)}
